\section*{HISTORICAL NOTE}
\phantomsection
\addcontentsline{toc}{section}{Historical Note}

(Numbers in brackets refer to the bibliography at the end of this note.)

The history of the theory of the multiplicative group \(R_{+}^{*}\) of real numbers \textgreater0 is closely related to that of the development of the notion of powers of a number \textgreater0 and the notations employed to denote them. The idea of the ``geometric progression'' formed by successive powers of the same number goes back to the Egyptians and the Babylonians; the Greek mathematicians were familiar with it, and as early as Euclid {[}1{]} we find a general statement equivalent to the rule \(a^{m}a^{n}=a^{m+n}\) for integral exponents m, n\textgreater0. In the Middle Ages, the French mathematician N. Oresme (14th century) rediscovered this rule. He was the first to have the idea of a fractional exponent \textgreater0, with a notation similar to our own and to our relevant rules of calculation, stated in general terms; for example, he used the two rules which we should write as

\[
(a b) ^ {1 / n} = a ^ {1 / n} b ^ {1 / n}, \quad (a ^ {m}) ^ {p / q} = (a ^ {m p}) ^ {1 / q} [ 2 ].
\]

But the ideas of Oresme were too far ahead of the mathematics of his time to exercise an influence on his contemporaries, and his treatise soon sank into oblivion. A century later, N. Chuquet stated Euclid's rule anew; furthermore, he introduced an exponential notation for the powers of unknowns in his equations, and did not hesitate to use the exponent o and integral exponents \(< 0\) (*). This time, even though Chuquet's work remained in manuscript and seems not to have been widely circulated, the idea of the isomorphism between the ``arithmetic progression'' of the exponents and the ``geometric progression'' of the powers was not lost sight of again; it was extended to negative and fractional exponents by Stifel {[}4{]}, and led finally to the definition of logarithms and the construction of the first tables, independently undertaken by the Scotsman J. Napier in 1614-1620 {[}5{]} and the Swiss J. Bürgi (whose work did not appear until 1620, although his ideas went back to the first years of the 17th century). Bürgi implicitly assumes the continuity of the isomorphism established between R and \(R^{*}_{+}\) by means of interpolation in the use of his tables; Napier on the other hand formulates it in his definition explicitly (or at any rate as explicitly as the hazy notion of continuity current at that time allowed) (*).

It is not our purpose here to dwell on the services rendered by logarithms to numerical calculation; from the theoretical point of view, their importance dates from the beginnings of the infinitesimal calculus, with the discovery of the expansions in series of \(\log(1+x)\) and \(e^{x}\) , and the differential properties of these functions. As far as the definition of logarithms and exponentials was concerned, it was assumed intuitively, up to the middle of the 19th century, that the function \(a^{x}\) defined for all rational x could be extended by continuity to the set of all real numbers; and it was not until the notion of real number had been definitively clarified and deduced from that of rational number that a rigorous justification was sought for this extension by continuity. An analogous principle of extension, suitably applied, is again at the base of the proofs of Propositions 1 and 2 of §2; from it there follows not only the definition of exponentials and logarithms, but also, as we shall see in Chapter VIII, angular measure.

